\begin{table}[t]
\centering
\caption{Descriptive statistics, estimation sample, 1976--2019}
\label{tab:descriptives}
\begin{threeparttable}
\begin{tabular}{lccc}
\toprule
 & \multicolumn{2}{c}{Adopting states} & Never-adopting \\
\cmidrule(lr){2-3}
 & Before adoption & After adoption & states \\
\midrule
Patents per year (fractional count) & 1,077 & 3,122 & 304 \\
 & (1,126) & (6,003) & (263) \\[2pt]
Patents per 100,000 residents & 20.4 & 35.7 & 11.4 \\
 & (17.4) & (22.7) & (5.9) \\[2pt]
Population (millions) & 5.6 & 7.1 & 2.6 \\
 & (5.1) & (7.9) & (1.8) \\[2pt]
Business R\&D, \$ millions (1997--2019) & 2,341 & 8,318 & 653 \\
 & (1,804) & (17,882) & (592) \\[2pt]
Share of patents (\%) held by: & & & \\
\quad small owners ($<$50 patents) & 20.2 & 20.0 & 24.4 \\
\quad medium owners (50--499) & 14.0 & 17.0 & 15.5 \\
\quad large owners (500+) & 36.9 & 50.2 & 32.8 \\
\quad no organization owner & 28.9 & 12.8 & 27.3 \\
\midrule
States & 23 & 23 & 9 \\
State-years & 446 & 566 & 396 \\
\bottomrule
\end{tabular}
\begin{tablenotes}[flushleft]\footnotesize
\item \textit{Notes:} Means across state-years, standard deviations in parentheses. Patents are
utility patents counted by the filing year and divided equally among inventors, so a patent with
inventors in two states adds one half to each. Owner size is the number of patents the first-listed
organization owner holds in the whole data. Business R\&D comes from NCSES and starts in 1997;
some state-years are suppressed. The before and after columns contain the same 23 states.
\end{tablenotes}
\end{threeparttable}
\end{table}
